\documentclass{zhvt-classic}

% 使用本地源流明体（GenRyuMin2 TW）按文件路径加载，不依赖系统字体名
% 版框（内框／界栏／书口侧栏／鱼尾）改用与版心印章同一印泥棕色：
%   印章图片 yshs.png 实色 = RGB(135,68,52) = HTML 874434
\usepackage{xcolor}
\usepackage{tikz}
% sealbrown 双注册：xcolor 供 TikZ/pgf，L3 供 cls 版框
\definecolor{sealbrown}{HTML}{874434}
\ExplSyntaxOn
\color_set:nnn { sealbrown } { HTML } { 874434 }
\ExplSyntaxOff

\zhvtset{
  grid color       = sealbrown,
  main font        = GenRyuMin2TW-SB,
  main font path   = C:/Users/Administrator/texmf/fonts/opentype/public/GenRyuMin2TW/,
jiazhu font      = GenRyuMin2TW-M,
  jiazhu font path = C:/Users/Administrator/texmf/fonts/opentype/public/GenRyuMin2TW/,
  % 页码框容纳 4 字：正常页「九十七」+前置空 1 格=4 行恰好；
  % 100+ 页「一百零一」共 5 行超出，末字「一/二/三」被框裁掉（2026-10-04 实测），故放宽到 5
  shukou page number chars = 5,
  % 书口书名框高从默认 6 字减到 4 字，把省出的 2 字高拨给 chapter 框，
  % 使 6 字书名+卷名放进 9 字高的 chapter 框后不再压线（2026-10-05）
  shukou title chars = 4,
}

\title{史记}[五宗世家]
\author{汉·司马迁}
\maker{闲叟漫记·重製}


% ===== 索隐述赞样式（2026-11-08 四次定稿 · 愚大人亲手改定参数）=====
% ⚠ 竖排前提：自定义字体家族必须带 RawFeature={vertical:+vert}；
%   simhei.ttf 加 +vhal 会整体加载失败静默回退，只写 +vert；simfang 同理。
\newCJKfontfamily{\shuzanfont}[Path=C:/Windows/Fonts/,RawFeature={vertical:+vert}]{simhei.ttf}
\newCJKfontfamily{\shuzanbody}[Path=C:/Windows/Fonts/,RawFeature={vertical:+vert}]{simfang.ttf}
% 题头：152x30pt 圆角 6pt 棕底白字；baseline (0,2) 使整框落入界行格子不压版框线；
% 白字节点 x 不补偿、y +8.2pt（像素级墨迹实测 dx/dy≈0，L/R 留白约 20pt）。
\newcommand{\shuzanhead}[1]{%
  \begin{tikzpicture}[baseline={(0,2)}]
    \fill[sealbrown, rounded corners=6pt]
      (-76pt,-15pt) rectangle (76pt,15pt);
    \node[text=white, inner sep=0pt, anchor=center] at (0pt,8.2pt)
      {\shuzanfont\fontsize{23}{23}\selectfont #1};
  \end{tikzpicture}%
}
% ★ 修正书口卷名位置：chapter 框起始位置改回 cls 默认值，避免全局影响。
%   只在每个卷首页通过 \juan 命令里的空命令微调该页位置（安全，不会跨页污染）。
\ExplSyntaxOn
\AtBeginDocument{
  \dim_set:Nn \l__zhvt_shukou_chapter_title_start_dim
    { \l__zhvt_shukou_title_height_dim + \l__zhvt_font_size_dim * 2 }
  \dim_set:Nn \l__zhvt_shukou_page_number_start_dim
    { \l__zhvt_shukou_chapter_title_start_dim + \l__zhvt_shukou_chapter_title_height_dim }
}
\ExplSyntaxOff

\begin{document}

% ★★ 版心（书口）排布覆写：把书名放到鱼尾之下（2026-10-04）
%
%   cls 默认版心自上而下：书名（title 静态框，在鱼尾之上）→ 鱼尾 → 篇名（chapter 动态框）→ 页码
%   本模板改为：          鱼尾（位置不动）→ 书名 → 篇名 → 页码分割线 → 页码
%
%   手法：把 cls 放在 title 静态框里的书名，改放进 chapter 动态框（鱼尾之下）；
%         title 静态框留空。书名取 \l__zhvt_title_str，即 \title{...} 的第一参数，
%         所以换书时只改 \title，不必在此重复写书名。
%
%   ⚠ 三条必须遵守的写法（2026-10-04 实测）：
%     ① 汉字必须用 \hbox{...} 包住。flowfram 会把框内容装进一个宽 14.4mm（= 一个 line_sep
%        = 书口栏宽）的纵向 minipage 并按该宽度折行，而整页靠 PDF /Rotate 90 显示、
%        框内文字本身不变向 —— 不包 \hbox 时「围炉夜话」4 字会被逐字折成 4 行，
%        在版面上就成了横排 4 字并横向溢出版框。
%     ② 竖排下**一个 \hbox 就是一整列**。书名与篇名要同处一列，必须放进同一个 \hbox；
%        分成两个 \hbox 时，后一个会被当成新的一列、甩到隔壁栏去
%        （实测「序」跑到书口栏外侧 14.4mm 的正文列上）。中间的 \quad 在框内是横向空白，
%        显示出来正好是书名与篇名之间的纵向间隔。
%     ③ ★ 内容盒要写成 **与框等宽的定宽盒** `\hbox to \l__zhvt_shukou_chapter_title_height_dim{…\hfil}`。
%        实测：chapter 框是把内容盒沿纵向**居中**摆放的，内容宽度一变，书名就跟着漂：
%        自然宽度的 \hbox（书名 38.4mm）落在 x=157.2mm，加篇名后（57.6mm）落在 x=152.4mm，
%        目录页（无篇名）与正文页的书名就差半格。做成等宽盒后，无论有没有篇名，
%        内容固定从框顶起排，实测两页均落在 x=142.8mm，不再漂。
\ExplSyntaxOn
% 把书口 chapter 框（书名+篇名）起始位置下移 1 字高，避免 4 字卷名顶到鱼尾
% ⚠ cls 的 \__zhvt_calculate_dimensions: 在 \AtBeginDocument 里执行，会重算该 dim；
%   必须把覆写挂到 begin document 之后，否则会被二次覆盖（2026-10-05 实测）。
% (旧方案已移到导言区，直接修改 chapter 框起始位置)
% 卷标题命令 #1=书口篇名(鱼尾) #2=正文卷名（含卷序，卷名内空格用 ~）
%
% ★★ 目录页两项修订（2026-10-05，愚大人指出）★★
%   ① 页码起点统一对齐：
%      原状 —— 目录项写成 \par\hspace{4em}卷名\hspace{4em}页码，页码起点 = 卷名实宽 + 4em。
%      卷名实宽随字数变：「卷一~梁惠王上」6 字 +~、「卷七~离娄上」5 字 +~，
%      差 1 个字位 → 卷七~卷十 4 条的页码整体比其余 10 条高 27.21pt（一个字），看着上下不齐。
%      改法 —— 卷名补全角空格（宽 = 1 字 = 27.215pt）补足到 **7 字符位**再写进目录，
%      14 条页码起点遂全部落在同一条竖线上（起点对齐；末字随页码字数自然参差）。
%      ⚠ 只对目录里的副本 \l_tmpc_tl 补，正文卷名仍用 \l_tmpa_tl 原样 —— 正文像素不受影响。
%   ② 卷名整体上移半字位：前导 \hspace 由 4em 改 3.5em。
%      实测原状目录块在版框内偏下（天头留白 139.6pt / 地脚留白 110.7pt），
%      上移 0.5em = 13.6pt 后上下基本均分，视觉居中。
\NewDocumentCommand { \juan } { mm } {
	\clearpage
	\tl_set:Nn \l__zhvt_chapter_title_tl { #1 }
	\tl_set:Nn \l_tmpa_tl { #2 }
	\tl_set_eq:NN \l_tmpc_tl \l_tmpa_tl
	\int_step_inline:nn { 7 - \str_count:n { #2 } }
	  { \tl_put_right:Nn \l_tmpc_tl { 　 } }
	\tl_set:Nx \l_tmpb_tl { \zhnum { page } }
	\tl_put_left:Nn \l_tmpb_tl { \hspace { 4em } }
	\tl_put_left:Nn \l_tmpb_tl { \l_tmpc_tl }
	\tl_put_left:Nn \l_tmpb_tl { \par \hspace { 3.5em } }
	\addtocontents { toc } { \l_tmpb_tl }
	\skip_horizontal:n { \l__zhvt_font_size_dim * 4 }
	\l_tmpa_tl
	\par
	\thispagestyle { empty }
}
\cs_set_protected:Npn \__zhvt_set_shukou_content: {
	% ★ 鱼尾上方留白（title 静态框）嵌入闲叟钤印：
	%   整页 /Rotate 90，故此处把图用 \rotatebox{90} 转正为竖章；
	%   尺寸卡在版心条幅内（书内约 7mm 宽 × 33.8mm 高，留余量不压鱼尾）。
	\setstaticcontents * { title } {
		\vbox_to_ht:nn { \l__zhvt_baseline_skip_dim } {
			\vskip 14.2pt
			\hbox_to_wd:nn { \l__zhvt_shukou_title_height_dim } {
				\hfil
				\rotatebox [ origin = c ] { 90 } { \includegraphics [ width = 7mm , height = 33.8mm ] { yshs.png } }
				\hfil
			}
		}
	}
% ★ 版心（书口）统一改法（2026-11-16 愚大人定）：18pt＋框内居中＋书口文字同版框色
%   ① 原 27.2pt 时 7 字长篇名（如「廉颇蔺相如列传」）末字溢出 chapter 框、压住页码；
%      改 \fontsize{18pt}{18pt}，前后 \hfil 令字串在框内居中（不随篇名字数漂移）。
%   ② \vbox{\vskip 8.2pt} 定量下移，抵消缩字号后栏宽方向 8.2pt 上偏（像素实测 y 中心 543.3 = 原版）。
%      ⚠ \raisebox / \vspace* 在 flowfram 框内无效；换字号须同步换补偿：16pt→10.1、19pt→7.4。
%   ③ chapter 与 page~number 均置 \color{sealbrown}（HTML 874434），与版框／界行同色；正文仍黑。
%   ⚠ 页码字位、钤印、鱼尾、界行、正文一律未动；五部合订本同此口径。
	\setdynamiccontents * { chapter }
	{ \color{sealbrown} \vbox { \vskip 8.2pt \hbox to \l__zhvt_shukou_chapter_title_height_dim
		{ \hfil \fontsize{18pt}{18pt}\selectfont \l__zhvt_title_str \quad \l__zhvt_chapter_title_tl \hfil } } }
	\setdynamiccontents * { page~number } { \color{sealbrown} \quad \zhnum { page } }
}
\__zhvt_set_shukou_content:
\ExplSyntaxOff

\maketitle

% 单篇，无多卷目录，不渲染目录页

\juan{五宗世家}{卷五十九~五宗世家第二十九}
\hspace*{2em}孝景皇帝子凡十三人为王。而母五人。同母者为宗亲。栗姬子曰荣。德。阏于。程姬子曰馀。非。端。贾夫人子曰彭祖。胜。唐姬子曰发。王夫人兒姁子曰越。寄。乘。舜。\\
\hspace*{2em}河间献王德。以孝景帝前二年用皇子为河间王。好儒学。被服造次必於儒者。山东诸儒多从之游。\\
\hspace*{2em}二十六年卒。子共王不害立。四年卒。子刚王基代立。十二年卒。子顷王授代立。\\
\hspace*{2em}临江哀王阏于。以孝景帝前二年用皇子为临江王。三年卒。无後。国除为郡。\\
\hspace*{2em}临江闵王荣。以孝景前四年为皇太子。四岁废。用故太子为临江王。\\
\hspace*{2em}四年。坐侵庙壖垣为宫。上徵荣。荣行。祖於江陵北门。既已上车。轴折车废。江陵父老流涕窃言曰。﹁吾王不反矣。﹂荣至。诣中尉府簿。中尉郅都责讯王。王恐。自杀。葬蓝田。燕数万衔土置冢上。百姓怜之。\\
\hspace*{2em}荣最长。死无後。国除。地入于汉。为南郡。\\
\hspace*{2em}右三国本王皆栗姬之子也。\\
\hspace*{2em}鲁共王馀。以孝景前二年用皇子为淮阳王。二年。吴楚反破後。以孝景前三年徙为鲁王。好治宫室苑囿狗马。季年好音。不喜辞辩。为人吃。\\
\hspace*{2em}二十六年卒。子光代为王。初好音舆马。晚节啬。惟恐不足於财。\\
\hspace*{2em}江都易王非。以孝景前二年用皇子为汝南王。吴楚反时。非年十五。有材力。上书原击吴。景帝赐非将军印。击吴。吴已破。二岁。徙为江都王。治吴故国。以军功赐天子旌旗。元光五年。匈奴大入汉为贼。非上书原击匈奴。上不许。非好气力。治宫观。招四方豪桀。骄奢甚。\\
\hspace*{2em}立二十六年卒。子建立为王。七年自杀。淮南。衡山谋反时。建颇闻其谋。自以为国近淮南。恐一日发。为所并。即阴作兵器。而时佩其父所赐将军印。载天子旗以出。易王死未葬。建有所说易王宠美人淖姬。夜使人迎与奸服舍中。及淮南事发。治党与颇及江都王建。建恐。因使人多持金钱。事绝其狱。而又信巫祝。使人祷祠妄言。建又尽与其姊弟奸。事既闻。汉公卿请捕治建。天子不忍。使大臣即讯王。王服所犯。遂自杀。国除。地入于汉。为广陵郡。\\
\hspace*{2em}胶西于王端。以孝景前三年吴楚七国反破後。端用皇子为胶西王。端为人贼戾。又阴痿。一近妇人。病之数月。而有爱幸少年为郎。为郎者顷之与後宫乱。端禽灭之。及杀其子母。数犯上法。汉公卿数请诛端。天子为兄弟之故不忍。而端所为滋甚。有司再请削其国。去太半。端心愠。遂为无訾省。府库坏漏尽。腐财物以巨万计。终不得收徙。令吏毋得收租赋。端皆去卫。封其宫门。从一门出游。数变名姓。为布衣。之他郡国。\\
\hspace*{2em}相。二千石往者。奉汉法以治。端辄求其罪告之。无罪者诈药杀之。所以设诈究变。彊足以距谏。智足以饰非。相。二千石从王治。则汉绳以法。故胶西小国。而所杀伤二千石甚众。\\
\hspace*{2em}立四十七年。卒。竟无男代後。国除。地入于汉。为胶西郡。\\
\hspace*{2em}右三国本王皆程姬之子也。\\
\hspace*{2em}赵王彭祖。以孝景前二年用皇子为广川王。赵王遂反破後。彭祖王广川。四年。徙为赵王。十五年。孝景帝崩。彭祖为人巧佞卑谄。足恭而心刻深。好法律。持诡辩以中人。彭祖多内宠姬及子孙。相。二千石欲奉汉法以治。则害於王家。是以每相。二千石至。彭祖衣皁布衣。自行迎。除二千石舍。多设疑事以作动之。得二千石失言。中忌讳。辄书之。二千石欲治者。则以此迫劫。不听。乃上书告。及汙以奸利事。彭祖立五十馀年。相。二千石无能满二岁。辄以罪去。大者死。小者刑。以故二千石莫敢治。而赵王擅权。使使即县为贾人榷会。入多於国经租税。以是赵王家多金钱。然所赐姬诸子。亦尽之矣。彭祖取故江都易王宠姬王建所盗与奸淖姬者为姬。甚爱之。\\
\hspace*{2em}彭祖不好治宫室。禨祥。好为吏事。上书原督国中盗贼。常夜从走卒行徼邯郸中。诸使过客以彭祖险陂。莫敢留邯郸。\\
\hspace*{2em}其太子丹与其女及同产姊奸。与其客江充有卻。充告丹。丹以故废。赵更立太子。\\
\hspace*{2em}中山靖王胜。以孝景前三年用皇子为中山王。十四年。孝景帝崩。胜为人乐酒好内。有子枝属百二十馀人。常与兄赵王相非。曰。﹁兄为王。专代吏治事。王者当日听音乐声色。﹂赵王亦非之。曰。﹁中山王徒日淫。不佐天子拊循百姓。何以称为籓臣。﹂\\
\hspace*{2em}立四十二年卒。子哀王昌立。一年卒。子昆侈代为中山王。\\
\hspace*{2em}右二国本王皆贾夫人之子也。\\
\hspace*{2em}长沙定王发。发之母唐姬。故程姬侍者。景帝召程姬。程姬有所辟。不原进。而饰侍者唐兒使夜进。上醉不知。以为程姬而幸之。遂有身。已乃觉非程姬也。及生子。因命曰发。以孝景前二年用皇子为长沙王。以其母微。无宠。故王卑湿贫国。\\
\hspace*{2em}立二十七年卒。子康王庸立。二十八年。卒。子鲋鮈立为长沙王。\\
\hspace*{2em}右一国本王唐姬之子也。\\
\hspace*{2em}广川惠王越。以孝景中二年用皇子为广川王。\\
\hspace*{2em}十二年卒。子齐立为王。齐有幸臣桑距。已而有罪。欲诛距。距亡。王因禽其宗族。距怨王。乃上书告王齐与同产奸。自是之後。王齐数上书告言汉公卿及幸臣所忠等。\\
\hspace*{2em}胶东康王寄。以孝景中二年用皇子为胶东王。二十八年卒。淮南王谋反时。寄微闻其事。私作楼车镞矢战守备。候淮南之起。及吏治淮南之事。辞出之。寄於上最亲。意伤之。发病而死。不敢置後。於是上。寄有长子者名贤。母无宠。少子名庆。母爱幸。寄常欲立之。为不次。因有过。遂无言。上怜之。乃以贤为胶东王奉康王嗣。而封庆於故衡山地。为六安王。\\
\hspace*{2em}胶东王贤立十四年卒。谥为哀王。子庆为王。\\
\hspace*{2em}六安王庆。以元狩二年用胶东康王子为六安王。\\
\hspace*{2em}清河哀王乘。以孝景中三年用皇子为清河王。十二年卒。无後。国除。地入于汉。为清河郡。\\
\hspace*{2em}常山宪王舜。以孝景中五年用皇子为常山王。舜最亲。景帝少子。骄怠多淫。数犯禁。上常宽释之。立三十二年卒。太子勃代立为王。\\
\hspace*{2em}初。宪王舜有所不爱姬生长男棁。棁以母无宠故。亦不得幸於王。王后脩生太子勃。王内多。所幸姬生子平。子商。王后希得幸。及宪王病甚。诸幸姬常侍病。故王后亦以妒媢不常侍病。辄归舍。医进药。太子勃不自尝药。又不宿留侍病。及王薨。王后。太子乃至。宪王雅不以长子棁为人数。及薨。又不分与财物。郎或说太子。王后。令诸子与长子棁共分财物。太子。王后不听。太子代立。又不收恤棁。棁怨王后。太子。汉使者视宪王丧。棁自言宪王病时。王后。太子不侍。及薨。六日出舍。太子勃私奸。饮酒。博戏。击筑。与女子载驰。环城过市。入牢视囚。天子遣大行骞验王后及问王勃。请逮勃所与奸诸证左。王又匿之。吏求捕勃大急。使人致击笞掠。擅出汉所疑囚者。有司请诛宪王后脩及王勃。上以脩素无行。使棁陷之罪。勃无良师傅。不忍诛。有司请废王后脩。徙王勃以家属处房陵。上许之。\\
\hspace*{2em}勃王数月。迁于房陵。国绝。月馀。天子为最亲。乃诏有司曰。﹁常山宪王蚤夭。后妾不和。適孽诬争。陷于不义以灭国。朕甚闵焉。其封宪王子平三万户。为真定王。封子商三万户。为泗水王。﹂\\
\hspace*{2em}真定王平。元鼎四年用常山宪王子为真定王。\\
\hspace*{2em}泗水思王商。以元鼎四年用常山宪王子为泗水王。十一年卒。子哀王安世立。十一年卒。无子。於是上怜泗水王绝。乃立安世弟贺为泗水王。\\
\hspace*{2em}右四国本王皆王夫人兒姁子也。其後汉益封其支子为六安王。泗水王二国。凡兒姁子孙。於今为六王。\\
\hspace*{2em}太史公曰。高祖时诸侯皆赋。得自除内史以下。汉独为置丞相。黄金印。诸侯自除御史。廷尉正。博士。拟於天子。自吴楚反後。五宗王世。汉为置二千石。去﹁丞相﹂曰﹁相﹂。银印。诸侯独得食租税。夺之权。其後诸侯贫者或乘牛车也。\\

% ===== 索隐述赞（司马贞《史记索隐》）· 统一样式 =====
% 题头「索隐述赞曰」五字并入反白框；正文每四句一列，定宽盒 \hfil 居列中；
% \fontsize{24}{41} 列距对齐全书界行（121+41k）；\clearpage 使述赞独占一页。
\clearpage
\shuzanhead{索隐述赞曰}\linebreak
{\shuzanbody\fontsize{24}{41}\selectfont
\hbox to \hsize{\hfil 景十三子。五宗亲睦。栗姬既废。临江折轴。\hfil}\linebreak
\hbox to \hsize{\hfil 阏于早薨。河间儒服。馀好宫苑。端事驰逐。\hfil}\linebreak
\hbox to \hsize{\hfil 江都有才。中山禔福。长沙地小。胶东造镞。\hfil}\linebreak
\hbox to \hsize{\hfil 仁贤者代。浡乱者族。兒姁四王。分封为六。\hfil}%
\par}
\end{document}
